%% LyX 2.5.1 created this file.  For more info, see https://www.lyx.org/.
%% Do not edit unless you really know what you are doing.
\documentclass[brazil]{beamer}
\usepackage[T1]{fontenc}
\usepackage{textcomp}
\usepackage[utf8]{luainputenc}
\setcounter{secnumdepth}{3}
\setcounter{tocdepth}{3}

\makeatletter
%%%%%%%%%%%%%%%%%%%%%%%%%%%%%% Textclass specific LaTeX commands.
% this default might be overridden by plain title style
\newcommand\makebeamertitle{\frame{\maketitle}}%
% (ERT) argument for the TOC
\AtBeginDocument{%
  \let\origtableofcontents=\tableofcontents
  \def\tableofcontents{\@ifnextchar[{\origtableofcontents}{\gobbletableofcontents}}
  \def\gobbletableofcontents#1{\origtableofcontents}
}

%%%%%%%%%%%%%%%%%%%%%%%%%%%%%% User specified LaTeX commands.


\PassOptionsToPackage{brazil}{babel}
\usepackage{babel}


\usetheme{Boadilla}
\usecolortheme{whale}

\title[Biblioteca Modular em C]{Desenvolvimento de Biblioteca Modular Avançada\\de Tópicos Matemáticos em C}
\subtitle{Estrutura, Compilação (Make), Bibliotecas Estática/Dinâmica, Testes e Git}

\author[Vlieger, Meira, Wallace]{
  André De Vlieger Lima \and
  Lucas Correa Meira \and
  Vitor Wallace de Moraes Alves
}

\institute[USP]{Instituto de Matemática e Estatística e Ciência da Computação}
\date{19 de Outubro de 2026}



\makeatother

\usepackage{babel}
\begin{document}
% --- SLIDE 1: CAPA ---

\begin{frame}
\titlepage 
\end{frame}
% --- SLIDE 3: PROBLEMA E DOMÍNIO ---

\begin{frame}

\frametitle{1. Descrição do Problema e Domínio}
\begin{block}{Objetivo do Projeto}
Criar um ecossistema coeso, modular e reutilizável de ferramentas
matemáticas e geométricas na linguagem C, voltado para cálculo computacional
e ambientes gráficos 
\end{block}
Domínio Abrangido (5 Módulos): 
\begin{itemize}
\item \textbf{vetores}: Operações vetoriais em 2D, 3D e ND; 
\item \textbf{geometria}: Cadarço (área), Ray-Casting e volumes 3D; 
\item \textbf{calculo}: Diferenças finitas, gradientes e integração numérica; 
\item \textbf{calculo\_numerico}: Raízes (Bissecção e Newton-Raphson); 
\item \textbf{polinomios}: Polinômios em listas encadeadas (qualquer grau). 
\end{itemize}
\end{frame}
% --- SLIDE 4: ARQUITETURA ---

\begin{frame}

\frametitle{2. Arquitetura do Software}

Princípios de Design: 
\begin{itemize}
\item Encapsulamento estrito e visibilidade mínima; 
\item Interfaces públicas isoladas na pasta "include/"; 
\item Implementação interna contida na pasta "src/". 
\end{itemize}
\begin{block}{Estrutura de Diretórios}
\texttt{projeto/}\\
 \texttt{~|– include/~~~~~\# Contratos da API (.h)}\\
 \texttt{~|– src/~~~~~~~~~\# Implementação interna (.c)}\\
 \texttt{~|– tests/~~~~~~~\# Testes automatizados}\\
 \texttt{~|– app/~~~~~~~~~\# Aplicação cliente (main.c)}\\
 \texttt{~|– build/~~~~~~~\# Objetos intermediarios (.o)}\\
 \texttt{~|– lib/~~~~~~~~~\# Bibliotecas geradas (.a e .so)}\\
 \texttt{~|– bin/~~~~~~~~~\# Executáveis finais}\\
 \texttt{~|– Makefile~~~~\# Automação de build}\\
 \texttt{~`– README.md~~~~\# Documentação e instruções} 
\end{block}
\end{frame}
% --- SLIDE 5: API PÚBLICA ---

\begin{frame}

\frametitle{3. Organização e Responsabilidade dos Módulos}
\begin{itemize}
\item \textbf{vetores} (\texttt{vetores.h / .c}): Operações em R3 e ND,
normas, produtos escalar e vetorial; 
\item \textbf{geometria} (\texttt{geometria.h / .c}): Algoritmo Shoelace,
Ray-Casting em polígonos, malhas 3D; 
\item \textbf{calculo} (\texttt{calculo.h / .c}): Diferenças finitas centradas,
gradiente, integrais simples e duplas; 
\item \textbf{calculo\_numerico} (\texttt{calculo\_numerico.h / .c}): Métodos
da Bissecção e Newton-Raphson; 
\item \textbf{polinomios} (\texttt{polinomios.h / .c}): Polinômios esparsos
com alocação dinâmica. 
\end{itemize}
\end{frame}
% --- SLIDE 6: APLICAÇÃO CLIENTE ---

\begin{frame}

\frametitle{4. Aplicação Cliente (app/main.c)}
\begin{block}{Respeito Estrito à Interface Pública}
A aplicação cliente atua como um consumidor independente da biblioteca. 
\end{block}
Regras de Encapsulamento Adotadas: 
\begin{itemize}
\item Inclusão exclusiva dos cabeçalhos (.h) presentes na pasta include/; 
\item Proibição total de incluir arquivos da pasta src/; 
\item Interação única e exclusiva por meio das funções da API pública; 
\item Nenhuma dependência de membros internos ou structs privadas. 
\end{itemize}
\end{frame}
% --- SLIDE 7: MAKEFILE ---

\begin{frame}

\frametitle{5. Processo de Compilação com Makefile}

Principais Alvos Padronizados: 
\begin{itemize}
\item \textbf{make all}: Compila os módulos, gera as bibliotecas e o cliente; 
\item \textbf{make clean}: Remove os artefatos de build/, lib/ e bin/; 
\item \textbf{make run}: Compila (se necessário) e executa a aplicação; 
\item \textbf{make test}: Compila e executa a suíte de testes; 
\item \textbf{make install}: Instala cabeçalhos e bibliotecas em PREFIX. 
\end{itemize}
\begin{block}{Compilação Incremental}
Cada objeto build/\%.o depende exclusivamente do seu fonte src/\%.c
e do seu cabeçalho include/\%.h, evitando recompilações desnecessárias. 
\end{block}
\end{frame}
% --- SLIDE 8: BIBLIOTECA ESTÁTICA ---

\begin{frame}

\frametitle{6. Geração e Utilização da Biblioteca Estática}
\begin{block}{Mecanismo da Biblioteca Estática (libapp.a)}
O utilitário 'ar' empacota os arquivos objeto (.o). O linker copia
o código diretamente para o binário final durante a compilação. 
\end{block}
Comandos de Terminal: 
\begin{exampleblock}{Comandos}
\texttt{\# 1. Compilação dos objetos}\\
 \texttt{gcc -Wall -Wextra -Iinclude -c src/{*}.c \&\& mv {*}.o build/}\\[0.1cm]
\texttt{\# 2. Empacotamento com 'ar'}\\
 \texttt{ar rcs lib/libapp.a build/{*}.o}\\[0.1cm] \texttt{\# 3. Compilação
e linkagem da Aplicação Cliente}\\
 \texttt{gcc -Iinclude app/main.c lib/libapp.a -lm -o bin/client\_estatico} 
\end{exampleblock}
\end{frame}
% --- SLIDE 9: BIBLIOTECA DINÂMICA ---

\begin{frame}

\frametitle{7. Geração e Utilização da Biblioteca Dinâmica}
\begin{block}{Mecanismo da Biblioteca Dinâmica (libapp.so)}
Exige código independente de posição (-fPIC). A resolução dos símbolos
ocorre em tempo de execução pelo carregador do sistema operacional. 
\end{block}
Comandos de Terminal: 
\begin{exampleblock}{Comandos}
\texttt{\# 1. Compilação com flag -fPIC}\\
 \texttt{gcc -Wall -Wextra -fPIC -Iinclude -c src/{*}.c}\\[0.1cm]
\texttt{\# 2. Geração do objeto compartilhado (.so)}\\
 \texttt{gcc -shared build/{*}.o -o lib/libapp.so}\\[0.1cm] \texttt{\#
3. Linkagem e Execucao}\\
 \texttt{gcc -Iinclude app/main.c -Llib -lapp -lm -o bin/client\_dinamico}\\
 \texttt{LD\_LIBRARY\_PATH=lib ./bin/client\_dinamico} 
\end{exampleblock}
\end{frame}
% --- SLIDE 12: CORREÇÃO DE BUG ---

\begin{frame}

\frametitle{10. Exemplo de Teste que Falhou e Correção}
\begin{block}{Falha Identificada no Gradiente}
No cálculo da derivada parcial em src/calculo.c, o teste esperava
2.0 para a função f(x,y)=x\texttwosuperior +y\texttwosuperior{} no
ponto (1,1), mas recebeu -2.0 (sinal invertido). 
\end{block}
Causa Raiz e Correção em src/calculo.c: 
\begin{alertblock}{Incorreto (Ordem Invertida)}
\texttt{return (l\_inferior - l\_superior) / (2.0 {*} h);} 
\end{alertblock}
\begin{exampleblock}{Corrigido (Ordem Correta: f(x+h) - f(x-h))}
\texttt{return (l\_superior - l\_inferior) / (2.0 {*} h);} 
\end{exampleblock}
Resultado: Após o ajuste, o comando 'make test' aprovou 100\% dos
módulos.
\end{frame}
% --- SLIDE 13: GIT E TAG V1.0 ---

\begin{frame}

\frametitle{11. Estratégia de Versionamento com Git}

Estratégia do Histórico: 
\begin{itemize}
\item Desenvolvimento modular em branches separadas para cada funcionalidade; 
\item Commits atômicos e com mensagens descritivas sobre o progresso; 
\item Integração do código na branch principal por meio de merges; 
\item Criação da tag anotada \textbf{v1.0} para delimitar a versão final
e estável do projeto. 
\end{itemize}
\end{frame}
% --- SLIDE 14: DIFICULDADES E DECISÕES ---

\begin{frame}

\frametitle{12. Dificuldades Encontradas e Decisões}
\begin{itemize}
\item \textbf{Resolução de Caminhos do .so}:\\
 Solução: Configuração do rpath (-Wl,-rpath=./lib) no Makefile para
facilidade no 'make run'. 
\item \textbf{Gestão de Memória em Listas e Vetores}:\\
 Solução: Responsabilidade explícita na API (cada função de criação
exige sua função correspondente de liberação). 
\item \textbf{Comparação de Ponto Flutuante}:\\
 Solução: Padronização do uso de margem de tolerância numérica nos
testes. 
\end{itemize}
\end{frame}
% --- SLIDE 15: CONCLUSÃO ---

\begin{frame}

\frametitle{13. Resultados Obtidos e Conclusão}

Resultados Finais: 
\begin{itemize}
\item 5 módulos matemáticos totalmente desenvolvidos e testados. 
\item Automação completa no Makefile (bibliotecas estática e dinâmica). 
\item Aplicação cliente demonstrando o uso correto da API pública. 
\item Histórico do Git estruturado com branches, merges e tag v1.0. 
\end{itemize}
\end{frame}

\end{document}
